\chapter*{Introduction générale}
\addcontentsline{toc}{chapter}{Introduction générale}

Dans un environnement économique caractérisé par une concurrence accrue et des fluctuations fréquentes des prix, les entreprises doivent s'appuyer sur des données fiables et actualisées afin de prendre des décisions stratégiques pertinentes.

Le groupe \textbf{CHO}, acteur majeur dans le secteur de l'huile d'olive, opère sur plusieurs marchés internationaux où le suivi des prix constitue un levier essentiel de compétitivité. En effet, la diversité des marchés, la multiplicité des concurrents et la variabilité des prix rendent la surveillance des données particulièrement complexe.

Actuellement, le processus de collecte et d'analyse des prix repose en grande partie sur des méthodes manuelles, notamment la collecte terrain et l'utilisation de fichiers Excel. Cette approche présente plusieurs limites, notamment en termes de fiabilité, de rapidité et d'exploitation des données.

Dans ce contexte, la problématique principale de ce projet est la suivante :

\begin{quote}
	extit{Comment automatiser et fiabiliser la collecte et l'analyse des données de prix afin d'améliorer la prise de décision ?}
\end{quote}

Pour répondre à cette problématique, ce projet vise à concevoir et développer une plateforme de \textbf{Business Intelligence} permettant :

\begin{itemize}[leftmargin=1.2cm]
	\item d'automatiser la collecte des données à partir de différentes sources ;
	\item de structurer et fiabiliser les données collectées ;
	\item de produire des indicateurs pertinents pour le pilotage ;
	\item de fournir des tableaux de bord interactifs facilitant l'analyse.
\end{itemize}

Ce projet s'inscrit dans une démarche de transformation des données en informations exploitables, dans le but d'améliorer la visibilité et la réactivité de l'entreprise face aux évolutions du marché.

Le présent mémoire est structuré comme suit :

\begin{itemize}[leftmargin=1.2cm]
	\item le premier chapitre présente le contexte de l'entreprise ainsi que l'étude de l'existant ;
	\item le deuxième chapitre traite l'analyse des besoins et la conception de la solution ;
	\item le troisième chapitre décrit la réalisation technique et les résultats obtenus.
\end{itemize}

\clearpage
